% SPDX-License-Identifier: CC-BY-SA-4.0
% Part: Lexical Analysis
% Section: Finite state automata
% Exercise: Definitions on automata

\begingroup

\begin{exercice}[Configurations et suites d'actions d'un automate]\label{exo:lexing/automata/definitions}

  On se donne l'automate $A$ suivant :
  \begin{tikzpicture}[automaton, baseline=(1.base), x=17mm, y=17mm]
    \state[initial]   (1) at (0,1) {$1$}; 
    \state            (2) at (1,1) {$2$}; 
    \state[accepting] (3) at (2,1) {$3$}; 
    \state            (4) at (0,0) {$4$}; 
    \state            (5) at (1,0) {$5$}; 
    \state[accepting] (6) at (2,0) {$6$}; 

    \path (1) edge[loop above] node       {$a$} (1);
    \path (1) edge             node[swap] {$a$} (4);
    \path (1) edge             node       {$b$} (2);
    \path (2) edge             node       {$a$} (3);
    \path (2) edge[loop above] node       {$b$} (2);
    \path (2) edge             node       {$b$} (6);
    \path (3) edge[loop above] node       {$a$} (3);
    \path (4) edge             node[swap] {$b$} (5);
    \path (4) edge             node       {$b$} (2);
    \path (5) edge             node[swap] {$b$} (2);
  \end{tikzpicture}

  \begin{question}
  \item Donner la définition mathématique de l'automate $A$ sous forme d'un tuple. 
  \end{question}
  
  Donner le graphe de la relation d'action entre toutes les configurations accessibles
  dans la reconnaissance de chacun des mots suivants et en déduire s'ils sont reconnus :
  \begin{question}
  \item $ab$
  \item $\varepsilon$
  \item $aabbaa$.
  \item Quel est le langage reconnu par l'automate $A$ ?
  \end{question}

\end{exercice}

\usetikzlibrary{decorations.pathmorphing}

\begin{correction}

  \begin{question}

  \item $\langle \Sigma, Q, I, F, \rightarrow \rangle$, avec $\Sigma = \{a, b\}$,  $Q = \{1, 2, 3, 4, 5, 6\}$,  $I = \{1\}$,  $F = \{3, 6\}$ et 
    $$\rightarrow = \left\{\begin{array}{lllll}
    \langle 1, a, 1\rangle, &
    \langle 1, a, 4\rangle, &
    \langle 1, b, 2\rangle, &
    \langle 2, a, 3\rangle, &
    \langle 2, b, 2\rangle, \\
    \langle 2, b, 6\rangle, &
    \langle 3, a, 3\rangle, &
    \langle 4, b, 5\rangle, &
    \langle 4, b, 2\rangle, &
    \langle 5, b, 2\rangle\,
  \end{array}\right\}$$

  \item Ni $2$ ni $4$ n'est un état final, donc $ab$ est rejeté par $A$. 
    \begin{center}
      \begin{tikzpicture}[x=15mm]
        \node (1) at (0,1) {$\langle ab, 1 \rangle$}; 
        \node (2) at (1,1) {$\langle a, 1 \rangle$}; 
        \node (3) at (1,0) {$\langle a, 4 \rangle$}; 
        \node (4) at (2,1) {$\langle \varepsilon, 2 \rangle$}; 
        \node (5) at (2,0) {$\langle \varepsilon, 5 \rangle$}; 

        \path (1) edge[leadsto] (2);
        \path (1) edge[leadsto] (3);
        \path (2) edge[leadsto] (4);
        \path (3) edge[leadsto] (4);
        \path (3) edge[leadsto] (5);
      \end{tikzpicture}
    \end{center}

  \item Mot $\varepsilon$. $\langle \varepsilon, 1\rangle$ est une configuration initiale et une configuration d'arrêt mais ce n'est pas une configuration finale.
    Donc $\varepsilon$ est rejeté par $A$.

  \item Mot $aabbaa$.  On a $\langle aabbaa, 1 \rangle \leadsto^\star \langle \varepsilon, 3 \rangle$, avec $3\in F$, donc $aabbaa$ est accepté. 
    \begin{center}
      \begin{tikzpicture}[x=20mm]
        \node (01) at (-.1,1) {$\langle aabbaa,      1 \rangle$}; 
        \node (11) at (1  ,1) {$\langle abbaa,       1 \rangle$}; 
        \node (10) at (1  ,0) {$\langle abbaa,       4 \rangle$}; 
        \node (21) at (2  ,1) {$\langle bbaa,        1 \rangle$}; 
        \node (20) at (2  ,0) {$\langle bbaa,        4 \rangle$}; 
        \node (31) at (3  ,1) {$\langle baa,         2 \rangle$}; 
        \node (30) at (3  ,0) {$\langle baa,         5 \rangle$}; 
        \node (41) at (4  ,1) {$\langle aa,          6 \rangle$}; 
        \node (40) at (4  ,0) {$\langle aa,          2 \rangle$}; 
        \node (50) at (5  ,0) {$\langle a,           3 \rangle$}; 
        \node (60) at (6  ,0) {$\langle \varepsilon, 3 \rangle$}; 

        \path (01) edge[leadsto] (11);
        \path (01) edge[leadsto] (10);
        \path (11) edge[leadsto] (21);
        \path (11) edge[leadsto] (20);
        \path (21) edge[leadsto] (31);
        \path (20) edge[leadsto] (31);
        \path (20) edge[leadsto] (30);
        \path (31) edge[leadsto] (41);
        \path (31) edge[leadsto] (40);
        \path (30) edge[leadsto] (40);
        \path (40) edge[leadsto] (50);
        \path (50) edge[leadsto] (60);
      \end{tikzpicture}
    \end{center}

  \item Le langage est $a^\star b^+ (a^+ \mid b)$.

  \end{question}

\end{correction}

\endgroup
\endinput
